\documentclass[11pt]{article}
\usepackage[utf8]{inputenc}
\usepackage{amsmath, amssymb}
\usepackage{geometry}
\geometry{a4paper, margin=1in}
\usepackage{hyperref}
\usepackage{ctex}

\title{Ullr's Secret 核心计算理论依据}
\author{}
\date{}

\begin{document}
\maketitle

\section*{概述}
本文档列出了 \texttt{core.py} 中所实现计算的具体物理公式与理论依据。主要面向程序员阅读，便于理解代码背后的数学模型。

\section{蒸汽压与湿度}
\subsection{\texttt{saturation\_vapor\_pressure(t\_celsius)}}
给定温度 $T$ ($^\circ\text{C}$)，其饱和蒸汽压 $e_s$ (hPa) 由马格努斯公式 (Magnus formula) 计算得出：
\[
e_s(T) = 6.112 \exp\left(\frac{17.67 T}{T + 243.5}\right)
\]

\subsection{\texttt{get\_dew\_point\_from\_rh(temp\_f, rh)}}
首先根据相对湿度 $RH$ (\%) 计算实际蒸汽压 $e$：
\[
e = e_s(T) \times \frac{RH}{100}
\]
然后通过马格努斯公式的逆运算求得露点温度 $T_d$ ($^\circ\text{C}$)：
\[
T_d = \frac{243.5 \ln(e / 6.112)}{17.67 - \ln(e / 6.112)}
\]

\subsection{\texttt{get\_rh\_from\_dew\_point(temp\_f, dew\_point\_f)}}
相对湿度通过露点 $T_d$ 的饱和蒸汽压与空气温度 $T$ 的饱和蒸汽压之比计算：
\[
RH = \left( \frac{e_s(T_d)}{e_s(T)} \right) \times 100\%
\]

\section{湿球温度}
\subsection{\texttt{psychrometric\_equation(tw\_c, t\_c, e, p\_hpa)} 与 \texttt{wet\_bulb\_f(...)}}
湿球温度 $T_w$ 的求解基于干湿球方程 (Psychrometric equation)：
\[
e_s(T_w) - P \cdot A \cdot (1 + 0.00115 T_w) \cdot (T - T_w) - e = 0
\]
其中 $P$ 为大气压 (hPa)，$A = 6.66 \times 10^{-4}$ 为干湿球常数。代码中使用 \texttt{fsolve} 求根算法找到满足该方程的 $T_w$。

\section{大气压力}
\subsection{\texttt{pressure\_at\_elevation(elevation\_ft)}}
给定海拔 $h$ (米) 的标准大气压 $P$ (hPa) 的近似公式为：
\[
P(h) = 1013.25 \left( 1 - 2.25577 \times 10^{-5} h \right)^{5.25588}
\]

\section{太阳位置}
\subsection{\texttt{estimate\_solar\_position(lat, lon, dt\_aware)}}
通过当前的 UTC 时间和经纬度，计算太阳的高度角 $\alpha$ 与方位角 $\gamma_s$。
\begin{align*}
b &= \frac{2\pi}{365}(N - 81) \quad \text{($N$ 为一年中的第几天)} \\
\delta &= 23.45^\circ \sin(b) \quad \text{(赤纬角)} \\
\text{EOT} &= 9.87 \sin(2b) - 7.53 \cos(b) - 1.5 \sin(b) \quad \text{(均时差)} \\
\text{LST} &= \text{UTC 时} + \frac{\text{经度}}{15} + \frac{\text{EOT}}{60} \quad \text{(地方太阳时)} \\
\omega &= 15^\circ (\text{LST} - 12) \quad \text{(时角)}
\end{align*}
高度角 $\alpha$：
\[
\sin(\alpha) = \sin(\phi)\sin(\delta) + \cos(\phi)\cos(\delta)\cos(\omega)
\]
方位角 $\gamma_s$：
\[
\cos(\gamma_s) = \frac{\sin(\delta) - \sin(\alpha)\sin(\phi)}{\cos(\alpha)\cos(\phi)}
\]

\section{辐射等效温度}
\subsection{\texttt{calculate\_radiative\_equivalent\_temps(...)}}
为评估太阳与热辐射对雪面温度的影响，使用传热系数 $K_{rad} = 8.0$ 计算短波辐射等效温度 $T_{SW,eq}$ 和长波辐射等效温度 $T_{LW,eq}$。

\subsubsection*{短波辐射等效 (Shortwave Equivalent)}
在坡度为 $\beta$、坡向为 $\gamma_a$ 的地形上，太阳入射角 $\theta_i$ 为：
\[
\cos(\theta_i) = \cos(\beta)\sin(\alpha) + \sin(\beta)\cos(\alpha)\cos(\gamma_s - \gamma_a)
\]
结合云量比例 $c$ 和地表反照率 $\alpha_s$，计算净短波辐射 $SW_{net}$：
\[
SW_{net} = 1361 \cdot \cos(\theta_i) \cdot (1 - 0.65 c^2) \cdot (1 - \alpha_s)
\]
\[
T_{SW,eq} = \frac{SW_{net}}{K_{rad}}
\]

\subsubsection*{长波辐射等效 (Longwave Equivalent)}
基于蒸汽压 $e_a$ 计算晴空发射率 $\epsilon_{clear}$：
\[
\epsilon_{clear} = 0.51 + 0.066 \sqrt{e_a}
\]
综合大气发射率 $\epsilon_{atm} = (1 - c)\epsilon_{clear} + c$。
考虑到地形遮挡，天空视角因子为 $SVF = \cos^2(\beta / 2)$，假定地形发射率 $\epsilon_{terrain} = 0.98$：
\[
LW_{net} = \sigma \left[ SVF \cdot \epsilon_{atm} T_{air}^4 + (1 - SVF) \cdot \epsilon_{terrain} T_{air}^4 - T_{snow}^4 \right]
\]
其中 $T_{snow} = 273.15$ K，斯特藩-玻尔兹曼常数 $\sigma = 5.67 \times 10^{-8}$。
\[
T_{LW,eq} = \frac{LW_{net}}{K_{rad}}
\]

\section{有效温度}
\subsection{\texttt{effective\_temperature\_f(...)}}
最终的雪面有效温度结合了热力学的湿球温度和辐射等效温度偏移：
\[
T_{eff} = T_w + T_{SW,eq} + T_{LW,eq}
\]

\section{雪层冻融物理模型 (Snow Consolidation Physics)}
\subsection{量纲转换}
模型中积分计算得出的有效积温单位为华氏度-小时 (F-hrs)。然而，标准的冰川水文学融化因子和物理结冰常数所适用的单位均为摄氏度-天 ($^\circ\text{C-days}$)。量纲转换常数为：
\[
1 \text{ } ^\circ\text{C-day} = 1.8 \times 24 = 43.2 \text{ F-hrs}
\]
因此，华氏度-小时积分 $I_{\text{F-hrs}}$ 转换为摄氏度-天积温 $I_{\text{C-days}}$ 的公式为：
\[
I_{\text{C-days}} = \frac{I_{\text{F-hrs}}}{43.2}
\]

\subsection{物理密度与热传导系数}
实际雪密度 $\rho$ 由雪水当量 (SWE, mm) 与实际物理雪深 $h_0$ (cm) 计算得出：
\[
\rho = \frac{\text{SWE}/10}{h_0}
\]
动态渗透融化系数 $K_M$ 与复冻热传导系数 $K_F$ 是基于真实密度的线性插值函数：
\begin{align*}
K_M &= 2.0\rho + 0.1 \\
K_F &= 10.0\rho + 0.5
\end{align*}

\subsection{融雪与复冻深度}
物理融雪深度 $D_{melt}$ (cm) 的计算基于度日融雪模型 (Degree-Day Method)：
\[
D_{melt} = K_M \cdot I_{m,\text{C-days}}
\]
其中 $I_{m,\text{C-days}}$ 为融化积分 ($T_{eff} > 32^\circ\text{F}$) 转换后的摄氏度-天。

物理复冻深度（即硬壳厚度）$D_{freeze}$ (cm) 的计算基于斯特凡相变方程 (Stefan Equation)：
\[
D_{freeze} = K_F \sqrt{I_{f,\text{C-days}}}
\]
其中 $I_{f,\text{C-days}}$ 为复冻积分绝对值 ($T_{eff} < 32^\circ\text{F}$) 转换后的摄氏度-天。

\end{document}